\documentclass[12pt]{article}
\usepackage{amsmath,amsfonts,amsthm,amssymb}
\usepackage{graphicx,float,wrapfig}
\usepackage{braket}

\topmargin=-0.45in      %
\evensidemargin=0in     %
\oddsidemargin=0in      %
\textwidth=6.5in        %
\textheight=9.0in       %
\headsep=0.25in         %

\linespread{1.1}

\title{\textbf{QCS 2026 Finals}}
\author{}
\date{}

\begin{document}
\maketitle

\section*{Intro}
The exam is on 2026.6.11 from 9:50 to 12:20. The overall score is 100. You may bring two A4 cheat sheets (double-sided).

\section{Concepts (20 pts)}
1. Name a discrete gate set for universal quantum computation. (2 pts)\\
2. List three different physical implementations of quantum computers. (3 pts)\\
3. Explain the concepts of complexity classes P, NP, BPP, and BQP. Give the relationship between P, BPP, and BQP. (5 pts)\\
4. For an error correction code encoding 1 logical qubit and being able to correct all single-qubit Pauli errors, what is the minimum number of physical qubits it can use? Show your reasoning and calculation. (5 pts)\\
5. Let $N=2^n$, and a function from $\{0,1\}^n$ to $\{0,1\}^m$ has period $r$, to which we apply a period finding quantum algorithm. Describe the characteristics of the outcome probability distribution after QFT, and write down the mathematical relation that the peak positions $y$ must satisfy. (5 pts)

\section{Gates and Physical Implementations (20 pts)}
1.\\
(a) We consider $qAND$ gate:
\[\ket a_1 \ket b_2=\ket{a \& b}_1\ket c_2, \]
where $a,b\in\{0,1\}$, and $\&$ is the classical AND operator. Can $\ket c_2$ be a qubit ancilla? Give your reason. (2 pts)\\
(b) We implement system $1,2$ with qutrits, whose computational basis are labeled $\ket 0, \ket 1, \ket 2.$ We use a gate called $iSWAP$ which swaps $\ket{02}$ and $\ket{11}$ and leaves all the other states as they are. Use this to construct the quantum circuit of $qAND$ and $qAND^{-1}$, and draw the circuit diagrams. (4 pts)\\
(c) How does this $qAND$ gate help to reduce circuit depth? (1 pts)\\
Hint: Consider gates with multiple control qubits, and in a physical implementation where the qutrits form some more complex structure than a $1-D$ chain but we can only perform $iSWAP$ on adjacent qutrits.\\ \\
2. We now consider a quantum computer using trapped ions. \\
(a) Explain the Doppler cooling mechanism for ions in a few sentences. (2 pts)\\
(b) Construct a circuit of $CNOT$ using Hadamard gates and $CZ$ gates, and draw the circuit diagram. (1 pts)\\
(c) Below we consider interaction between two of the multiple ions in the system. When considering ion $m$ and ion $n$, we can write their state as $\ket {e_i}_m \ket{e_j}_n\ket{s}. $ $\ket {e_i}$ can take the states $\ket g,\ket{e_0},\ket{e_1}$, corresponding the the ground state, the first excited state and the second excited state, respectively. The number $s$ stands for the shared sonic state number, which can be viewed as the energy level state for a quantum harmonic oscillator. The $\hat a$ and $\hat a^\dagger$ below stand for the annihilation operator and creation operator for such oscillator. When we are considering only one ion $m$ and the gate acts trivially on the $\ket {e_i}$ for the other ion, we can simply write the state we are focusing on as $\ket {e_i}_m \ket{s}$ for convenience. \\
Define a unitary operator
\[\hat U_{n}^{k,q}=\exp{[e^{-ik\pi /2}(\ket{e_q}_n\bra{g}\hat ae^{-i\phi}+\hat a^{\dagger}\ket{g}_n\bra{e_q}e^{i\phi})]}.\]
$k$ is a angle parameter of the gate. 
Calculate the state after applying the unitary on the following states: $\ket g_n\ket 0, \ket g_n\ket 1, \ket {e_q}_n\ket 0.$ (6 pts)\\
Hint: Use the Taylor expansion
\[e^{\hat A}=I+\sum_{n\geq 1}\frac{A^n}{n!}.\]
(d) Prove that if we encode on each ion the computational basis $\ket{0_L}=\ket{g}, \ket{1_L}=\ket{e_0}$, prove that the unitary
\[\hat U_{mn}=\hat U_m^{1,0}\hat U_n^{2,1}\hat U_m^{1,0}\]
is equivalent to the $CZ$ gate (4 pts).\\

\section{Phase matched Grover algorithm (20 pts)}
The Grover search algorithm searches for elements in a specific set $S$, containing $M$ elements, in the whole set containing $N>M$ elements.\\
1. Consider
\[U_f=I-2\sum_{x\in S}\ket x\bra x, D=2\ket s\bra s-I.\]
\[\ket s=\frac{1}{\sqrt N}\sum_{x=0}^{N-1}\ket x.\]
(a) Prove that $U_f$ is unitary. Write down the effect of $U_f$ on the basis states that are in $S$ and on those that are not. (3 pts + 2 pts)\\
(b) Calculate $DU_f\ket s$, and find out the relation between $M$ and $N$ such that the algorithm find a vector in $S$ with probability 1. (5 pts)\\
2. Consider
\[U_f=I-(1-e^{i\phi})\sum_{x\in S}\ket x\bra x, D=(1-e^{i\phi})\ket s\bra s-I.\]
(a) Prove that $U_f$ is unitary. Write down the effect of $U_f$ on the basis states that are in $S$ and on those that are not. (3 pts + 2 pts)\\
(b) Calculate $DU_f\ket s$, when $M/N$ exceeds a certain limit, the algorithm find a vector in $S$ with probability 1. Find the limit. (5 pts)\\
Hint: Expand using
\[\ket m=\frac{1}{\sqrt M}\sum_{x\in S}\ket x,\ket u=\frac{1}{\sqrt {N-M}}\sum_{x\notin S}\ket x.\]

\section{Deutsch-Jozsa Algorithm (20 pts)}
We want to determine whether a function $f:\{0,1\}^n\xrightarrow{} \{0,1\}$ is balanced or constant using the oracle
\[U_f\ket x \ket y=\ket x\ket{y\oplus f(x)}.\]
1. Prove that (6 pts)
\[U_f\ket x\ket -=(-1)^{f(x)}\ket x \ket -.\]
2. If we prepare the state to be $\ket 0^{\otimes n}\ket 1$, apply $H^{\otimes(n+1)}$, then $U_f$, then $H^{\otimes n}\otimes I$. Prove that if $f$ is constant then we always get $\ket 0^{\otimes n}$, but if it is balanced we never get $\ket 0^{\otimes n}$. (10 pts)\\
3. Consider the function
\[f(x_0,x_1)=x_1 \oplus x_0.\]
Calculate the final state in the Deutsch-Jozsa algorithm, and determine whether the function is balanced or constant. Does it fit with the result of direct counting? (4 pts)

\section{[[4,2,2]] Iceberg Code (20 pts)}
We consider the stabilizer code, encoding 2 logical qubits, with stabilizers
\[g_1=X_1X_2X_3X_4, g_2=Z_1Z_2Z_3Z_4.\]
1. Verify that $g_1,g_2$ commute and are independent of each other. (4 pts)\\
2. We define the logical $X$ and $Z$ operators as
\[\bar X_1=X_1X_2I_3I_4,\bar X_2=I_1X_2I_3X_4,\bar Z_1=Z_1I_2Z_3I_4,\bar Z_2=I_1I_2Z_3Z_4.\]
(i) Verify that all of those operators commute with $g_1$ and $g_2$. (2 pts)\\
(ii) Verify that all of those operators are not generated ny $g_1$ and $g_2$. (2 pts)\\
(iii) Verify that except for $\bar X_i \bar Z_i=-\bar Z_i\bar X_i,i=1,2$, all other pairs of operators in $\bar X_1,\bar X_2,\bar Z_1,\bar Z_2$ commute. (2 pts)\\
3. Find out the explicit computational basis expansion of the logical states
\[\ket{\overline{00}},\ket{\overline{01}},\ket{\overline{10}},\ket{\overline{11}}.\]
They satisfy $\bar Z_i\ket{\overline{a_1a_2}}=(-1)^{a_i}\ket{\overline{a_1a_2}},i=1,2.$ (4 pts)\\
4. Fill in the blanks with $+1$ and $-1$ that represent the error syndromes that can be obtained by measuring the stabilizers after the error occurred, reflecting whether the state is in the $+1$ or $-1$ eigenstate of the stabilizers.
\begin{table*}[ht]
\centering
\begin{tabular}{|c|*{4}{c|}|*{4}{c|}|*{4}{c|}}   % 1 Error, 4 X's, 3 Z's, 2 Y's
\hline
Error & X$_1$ & X$_2$ & X$_3$ & X$_4$ & Z$_1$ & Z$_2$ & Z$_3$ & Z$_4$ & Y$_1$ & Y$_2$ & Y$_3$ & Y$_4$ \\
\hline
S$_1$ &       &       &       &       &       &       &       &       &       &       &       &       \\
S$_2$ &       &       &       &       &       &       &       &       &       &       &       &       \\
\hline
\end{tabular}
\end{table*}
What observation can you make on those error syndromes for same-type errors on different qubits? How is it consistent with the code distance $d=2$? Why can the code detect all the single-qubit errors but cannot correct them? (6 pts)

\end{document}